\documentclass[11pt,letterpaper]{article}
\usepackage[T1]{fontenc}
\usepackage[utf8]{inputenc}
\usepackage{mathptmx}
\usepackage[scaled=0.92]{helvet}
\usepackage[margin=1in]{geometry}
\usepackage{graphicx}
\usepackage{booktabs}
\usepackage{tabularx}
\usepackage{array}
\usepackage{multirow}
\usepackage{amsmath}
\usepackage{xcolor}
\usepackage{listings}
\usepackage{caption}
\usepackage{enumitem}
\usepackage{xurl}
\usepackage[hidelinks]{hyperref}
\usepackage{textcomp}

\definecolor{cardinal}{HTML}{8C1515}
\definecolor{codebg}{HTML}{F7F6F3}
\hypersetup{colorlinks=true,linkcolor=cardinal,citecolor=cardinal,urlcolor=cardinal}
\captionsetup{font=small,labelfont=bf}
\setlist{itemsep=2pt,topsep=3pt}
\lstset{basicstyle=\ttfamily\scriptsize,breaklines=true,breakatwhitespace=true,columns=fullflexible,
  keepspaces=true,frame=single,rulecolor=\color{gray!40},backgroundcolor=\color{codebg},
  xleftmargin=4pt,xrightmargin=4pt,aboveskip=6pt,belowskip=6pt,showstringspaces=false}
\newcommand{\coe}{kgCO$_2$e/m$^2$}
\newcommand{\figdir}{../figures}

\title{\textbf{Can an LLM Estimate the Embodied Carbon of a Building Enclosure?}\\[2pt]
\large Error Analysis and Prompt Design for Early-Stage A1--A3 Estimates on 292 North American Buildings}
\author{George Zheng \and Zixuan Zhao \and Heshan Chen \and Fei Cheng \and Jamie Liu\\[4pt]
\small CEE226, Stanford University --- Group 8\\
\small Part of the project \emph{Climate- and Use-Specific Embodied Impacts of Exterior Wall Design:}\\
\small \emph{A Comparative LCA and Implications for LEED v5 Embodied Carbon Credits}}
\date{November 2026}

\emergencystretch=2em
\begin{document}
\maketitle

\begin{abstract}
\noindent We test whether a large language model (Claude Opus~5.5, run as a Claude Code subagent with medium reasoning effort) can estimate the cradle-to-gate (A1--A3) embodied carbon of the building enclosure from early-stage design features. The reference values come from the 292 buildings of the harmonized North American WBLCA dataset of Benke et al.~\cite{benke2025}, normalized per m$^2$ of floor area. With a plain prompt, only 19.5\% of estimates fall within $\pm$20\% of the reference LCA. An error analysis shows that the model compresses all buildings toward a typical value (log--log slope 0.46), never considers HFC-blown XPS insulation (15.5\% of all enclosure carbon, and about half of the systematic underestimate), underestimates envelope area when it is not given, and misjudges renovations. Its largest overestimates trace back to incomplete reference LCAs. A first revised prompt that added qualitative material knowledge and a step-by-step method made the model more consistent but \emph{less} accurate: the model repeated the knowledge without applying it to the numbers. A second revision placed externally sourced quantitative anchors (EPD and CLF baseline intensities) at the step where each number is chosen. It raised coverage to 25.0\%, lowered the bias from 0.75 to 0.86 (estimate/reference, geometric mean), and lowered the run-to-run CV from 19.0\% to 12.5\%. On a held-out test set of 88 buildings the gains point the same way but are not statistically significant. A feature study shows that ten early-design inputs are as informative as all 48; the binding constraint is missing assembly information, not prompt length or feature count. All data, prompts, code and model outputs are public.\footnote{\url{https://github.com/HeshanChen/CEE226LCA}}
\end{abstract}

\section{Introduction}
Early design decisions fix most of a building's embodied carbon, and the exterior enclosure is one of the most consequential and variable of them. Early-stage LCA is hard because design information is sparse, and the usual hand-off from architect to LCA practitioner is slow. Recent work shows that LLM in-context learning can support early-stage whole-building LCA, with domain knowledge and ontologies improving on innate knowledge~\cite{jung2026}. We narrow that question to the enclosure and to the product stage (A1--A3), where material choices dominate and where LEED~v5 embodied-carbon credits are earned in practice. We ask:
\begin{enumerate}[label=\textbf{Q\arabic*}]
\item How accurately and consistently does an LLM estimate enclosure A1--A3 embodied carbon from early-stage features, and why does it fail?
\item Which prompt changes improve the estimates, and which do not?
\item Which input features does the model actually need?
\end{enumerate}

\section{Data and Ground Truth}
\textbf{Source.} The harmonized dataset of Benke et al.~\cite{benke2025} (CC~BY~4.0) contains 927,950 material-level LCA results for 292 projects in the US and Canada, with 48 early-stage features per project (49 columns including the project index).

\textbf{Scope and unit.} We keep OmniClass element \texttt{Shell - Enclosure} and stage \texttt{A1-A3} (55,807 rows) and divide the total by floor area. We use \texttt{bldg\_cfa} for new construction and \texttt{bldg\_renovated\_gfa + bldg\_added\_gfa} for the 49 renovations, which reproduces the dataset's own \texttt{eci\_a\_to\_c\_cfa} normalization for all 292 projects. The enclosure includes windows, doors and roofs, and the unit is per m$^2$ of \emph{floor} area. This differs from the proposal's functional unit (1~m$^2$ of wall), because wall area is missing for most projects.

\textbf{Distribution.} The reference values have a median of 78.0~\coe{}, a mean of 95.6 and a range of 0.49--513.9 (5th--95th percentile 10.5--255.7). The enclosure is a median 19\% of whole-building A--C carbon.

\begin{figure}[t]
\centering
\includegraphics[width=0.62\linewidth]{\figdir/fig_composition.pdf}
\caption{Share of enclosure A1--A3 GWP by material group across the 292 projects. XPS alone is 15.5\%, at 43.6~kgCO$_2$e/kg, against 1.2 for mineral wool and 3.1 for polyiso. This is consistent with the HFC-blown XPS modelled by Tally and One Click LCA in the assessment years. 149 projects contain XPS.}
\label{fig:composition}
\end{figure}

\textbf{Composition.} Insulation (26.4\%) and aluminum (22.0\%) dominate the enclosure total (Fig.~\ref{fig:composition}).

\textbf{Data quality.} Ten new-construction projects have no window, glazing, curtain-wall or storefront GWP at all, although several report a window-to-wall ratio of 0.5--0.7. Project 74, for example, has WWR 0.65 and a reference value of 0.64~\coe{}. We keep these projects in every headline metric and also report results without them (``excl.\ incomplete'').

\section{Methods}
\textbf{Model and isolation.} Each estimate is produced by a Claude Code subagent (Claude Opus~5.5, reasoning effort medium) that may read only its task file and write only its answer files. It has no web access, no code execution and no access to the ground truth. A task file contains the system prompt, a shared context (the feature glossary) and 48--49 building blocks, each with the building's non-empty features followed by the task text. Appendix~\ref{app:prompts} gives the exact text. The model is told to treat every building as an independent request.

\textbf{Protocol.} Each prompt version is run 3 times on all 292 buildings (876 estimates). For every run, the buildings are shuffled with a fixed seed ($1000 + $ run index) and split into six chunks. The chunks are identical across prompt versions, so results can be paired building by building. The estimate for a building is the mean of the midpoints of its returned ranges. In one v3 run a container restart interrupted the subagents after 130 answers; the remaining 162 buildings were answered by continuation tasks with the same prompt.

\textbf{Parsing.} We use the course parser (the last line \texttt{RANGE\_kgCO2e\_per\_m2: low - high}) with one change: a lower bound of 0 is accepted. Renovations that leave the envelope untouched have reference values near zero, and the original parser would have discarded exactly those answers.

\textbf{Metrics.}
\begin{itemize}
\item Accuracy: the share of buildings whose reference value lies within the estimate $\pm$20\% (coverage, with Wilson 95\% CI), RMSE, median absolute percentage error (APE), and the median $|\log(\text{estimate}/\text{reference})|$.
\item Bias: the geometric mean of estimate/reference.
\item Consistency: the mean coefficient of variation (CV) across runs.
\item Paired comparisons between prompt versions: McNemar's test on coverage and the Wilcoxon signed-rank test on $|\log(\text{estimate}/\text{reference})|$.
\end{itemize}

\textbf{Held-out split.} No independent dataset with the same structure exists (Section~\ref{sec:heldout}). We therefore split the 292 projects into a dev set of 204 and a test set of 88 (seed 226), stratified on input features only: new construction vs.\ renovation $\times$ climate group.

\section{Baseline and Error Analysis}
\label{sec:error}
With prompt v1 (Appendix~\ref{app:prompts}), all 876 answers parsed.
\begin{itemize}
\item Coverage is 19.5\% (CI 15.4--24.5\%), RMSE 79.7~\coe{}, median APE 37.9\%, bias 0.75 and CV 19.0\% (Table~\ref{tab:main}).
\item 63\% of buildings are underestimated beyond the $\pm$20\% band, and 17.5\% overestimated.
\end{itemize}

\begin{figure}[t]
\centering
\includegraphics[width=0.49\linewidth]{\figdir/fig_bins.pdf}\hfill
\includegraphics[width=0.49\linewidth]{\figdir/fig_diagnostics.pdf}
\caption{Left: median estimate/reference by reference-value bin. All prompts compress toward a typical value; v3 shifts the curve up. Right: median estimate/reference for the subgroups behind each failure mode.}
\label{fig:diag}
\end{figure}

\textbf{Compression, not uniform underestimation.}
\begin{itemize}
\item The log--log slope of estimate on reference is 0.46: buildings below 25~\coe{} are overestimated 1.58$\times$, and buildings above 200 are estimated at 0.28$\times$ (Fig.~\ref{fig:diag}).
\item The estimates span 1.7--168~\coe{}, against 0.5--514 for the reference.
\item The 20 largest errors, all underestimates, contribute 68\% of the squared error.
\item The model leans on geometry. The envelope-to-floor ratio correlates 0.62 with the log estimate but only 0.39 with the log reference. What it cannot predict is the carbon intensity per m$^2$ of envelope, which correlates $-0.79$ with the log error. That intensity is set by assembly and material choices, and \emph{none of the 48 features describes the wall assembly}.
\end{itemize}

\textbf{Why the model underestimates.}
\begin{enumerate}
\item \emph{HFC-blown XPS is never considered.} XPS appears in 149 projects, yet only 0.6\% of answers mention it and none mention blowing agents. Median estimate/reference is 0.67 with XPS and 0.80 without. Removing XPS from the reference raises the median ratio from 0.72 to 0.84, so XPS explains about half of the bias.
\item \emph{Climate is an insulation effect.} The ratio is 0.65 in cold zones (5--8) and 0.77 in warm zones (1--3). XPS is present in 65--68\% of zone 5A/5B projects but only 26--33\% of zone 3B/3C projects.
\item \emph{Missing envelope area.} The ratio is 0.81 when \texttt{bldg\_therm\_env\_area} is given and 0.65 when it is not; the model's own take-off is too small.
\item \emph{Curtain-wall tail.} Single projects reach 300~\coe{} from aluminum curtain wall alone.
\item \emph{Renovations scoped too narrowly.} 60\% of major renovations are underestimated, and tenant improvements reach 0.24.
\item \emph{Scope wording.} The v1 scope text was narrower than the dataset's enclosure. For example, 11\% of enclosure GWP was tagged ``Superstructure'' in the source tool, mostly roof insulation and membranes.
\end{enumerate}

\textbf{Why the model overestimates.} The ten flagged incomplete references have a median ratio of 3.09; here the model's estimate is a plausible building and the reference is not. The remaining overestimates are mostly tenant improvements and minor renovations, where the features do not reveal how much envelope was replaced.

\textbf{Ruled out.} R-values, WWR, size, completion year and the embodied-carbon-reduction flag each correlate weakly with the error ($|r| \le 0.26$).

\section{Prompt Iterations}
\subsection{Prompt v2: qualitative knowledge and a step-by-step method}
Prompt v2 addressed each finding:
\begin{itemize}
\item a scope written to match the dataset's enclosure;
\item a required component take-off (areas $\rightarrow$ assemblies $\rightarrow$ intensities $\rightarrow$ total);
\item named high-impact materials: HFC-blown XPS at ``roughly ten times or more'' the GWP of other insulation per unit of thermal resistance~\cite{alberta,vermont}, and aluminum curtain wall at 100--200~\coe{} of facade~\cite{oneclick,arup};
\item an instruction not to anchor on a typical building;
\item guidance for renovations;
\item a \texttt{BEST} line in addition to the range.
\end{itemize}

\begin{table}[t]
\centering\small
\caption{Accuracy and consistency by prompt version (3 runs each). Coverage: reference within estimate $\pm$20\%. Bias: geometric mean of estimate/reference.}
\label{tab:main}
\begin{tabular}{llcccccc}
\toprule
Buildings & Prompt & Coverage (95\% CI) & RMSE & Median APE & Med.\ $|\log|$ & Bias & CV \\
\midrule
\multirow{3}{*}{All (292)} & v1 & 19.5\% (15.4--24.5) & 79.7 & \textbf{37.9\%} & 0.456 & 0.75 & 19.0\% \\
 & v2 & 15.1\% (11.4--19.6) & 81.0 & 44.3\% & 0.563 & 0.67 & 13.8\% \\
 & v3 & \textbf{25.0\%} (20.4--30.3) & \textbf{76.0} & 40.0\% & \textbf{0.426} & \textbf{0.86} & \textbf{12.5\%} \\
\midrule
\multirow{2}{*}{New constr.\ (243)} & v1 & 18.9\% & 77.4 & 37.5\% & -- & 0.76 & 17.3\% \\
 & v3 & \textbf{26.3\%} & \textbf{74.1} & 38.0\% & -- & \textbf{0.91} & \textbf{10.9\%} \\
\midrule
\multirow{3}{*}{Test (88)} & v1 & 20.5\% (13.3--30.0) & 74.1 & \textbf{36.4\%} & \textbf{0.393} & 0.75 & 19.2\% \\
 & v2 & 18.2\% (11.5--27.5) & 75.7 & 43.4\% & 0.511 & 0.66 & 15.2\% \\
 & v3 & \textbf{26.1\%} (18.1--36.2) & \textbf{69.2} & 37.0\% & 0.400 & \textbf{0.86} & \textbf{11.2\%} \\
\bottomrule
\end{tabular}
\end{table}

\textbf{Result: more consistent, less accurate.}
\begin{itemize}
\item v2 lowered the CV from 19.0\% to 13.8\% and raised the log--log slope to 0.55.
\item But coverage fell to 15.1\% and the bias to 0.67 (Table~\ref{tab:main}). The errors were larger (Wilcoxon $p<0.0001$), and v2 was closer to the reference for only 34\% of buildings.
\end{itemize}
Three mechanisms explain this:
\begin{enumerate}
\item \emph{Knowing is not applying.} v2 answers mention XPS 50 times more often, but 89\% still assume polyiso roof insulation and place XPS below grade. The median roof intensity stays at 35~\coe{}, and projects with XPS fall from 0.67 to 0.58.
\item \emph{Low bottom-up intensities.} Opaque walls were taken at 80, roofs at 35 and curtain wall at 145~\coe{}, giving 66.8 per m$^2$ of envelope against 86.7 in the reference. The prompt's curtain-wall range became the default value.
\item \emph{A rule that subtracts.} ``Count only new materials'' pushed minor renovations from 0.85 to 0.56.
\end{enumerate}
The self-reported \texttt{BEST} value sat below the range midpoint and was less accurate.

\subsection{Prompt v3: quantitative anchors at the decision step}
v3 keeps the v2 scope, method and anti-anchoring rule and changes four things (Appendix~\ref{app:v3}):
\begin{enumerate}
\item A reference table of A1--A3 intensities per m$^2$ of component, compiled by a separate research subagent from external sources, explicitly excluding the ground-truth dataset and the CLF WBLCA benchmark. The sources are CLF North American Material Baselines 2023/2025, PCI precast EPDs, Kawneer fenestration EPDs, RDH wall studies and others~\cite{clf2025,pci,kawneer,rdh}. The table covers wall types, fenestration, roof assemblies by insulation type (HFC XPS: 85--110 counting the product stage only, up to 335--410 when blowing-agent emissions are booked with the product), insulation per m$^2\cdot$RSI and aluminum per kg.
\item The tool modelling convention: LCA models of this period often model roof and below-grade insulation as HFC XPS, so the range must include that case when the insulation type is unknown.
\item A per-envelope sanity check, with a threshold of about 50~\coe{} of envelope derived from the component ranges.
\item A renovation rule that no longer subtracts, and a single \texttt{RANGE} output.
\end{enumerate}

\begin{figure}[t]
\centering
\includegraphics[width=0.78\linewidth]{\figdir/fig_scatter_v1_v3.pdf}\\[4pt]
\includegraphics[width=0.85\linewidth]{\figdir/fig_prompts.pdf}
\caption{Top: estimate vs.\ reference for v1 and v3 (line 1:1, band $\pm$20\%). Bottom: coverage with 95\% CI and median $|\log(\text{estimate}/\text{reference})|$ for all 292 buildings and the held-out 88.}
\label{fig:prompts}
\end{figure}

\textbf{Result.}
\begin{itemize}
\item v3 is the best prompt on coverage (25.0\%), RMSE (76.0), log error (0.426), bias (0.86) and consistency (CV 12.5\%) (Table~\ref{tab:main}, Fig.~\ref{fig:prompts}).
\item Against v2 the gain is significant (McNemar $p=0.002$; Wilcoxon $p<0.0001$). Against v1 it is consistent but borderline (McNemar $p=0.085$; Wilcoxon $p=0.053$; v3 closer for 57\% of buildings).
\item 96\% of answers for XPS projects now consider HFC XPS.
\item The non-XPS part of the enclosure is unbiased: estimate/(reference $-$ XPS) went from 0.84 to 0.99. Buildings with a given envelope area are also unbiased (0.99), and warm zones reach 0.98.
\item What remains: compression at the top (0.38 above 200~\coe{}); more overestimation of light enclosures (1.31 in the 25--50 bin); 0.73 when the envelope area is missing; and renovations (0.71--0.74).
\item The median APE is slightly above v1 (40.0\% vs.\ 37.9\%), because v3 shifts estimates up and APE penalizes overestimates more than underestimates.
\end{itemize}

\section{Held-out Evaluation}
\label{sec:heldout}
\textbf{No external dataset.} The CLF WBLCA Benchmark v2 probably shares source models with the ground truth. The open CarbEnMats database~\cite{carbenmats} reports envelope A1--A3 for only 172 European buildings, uses a different element definition, and has no geometry fields for them. We therefore hold out 88 test buildings.

\textbf{Post hoc split.} The split was assigned after the runs, and the error analysis had seen aggregate patterns across all 292 buildings. To check that this did not drive the prompt design, we re-ran every diagnostic on the 204 dev buildings alone. Every finding replicates, including:
\begin{itemize}
\item the log--log slope (0.44 vs.\ 0.46);
\item the XPS share (14.7\% vs.\ 15.5\%) and the XPS effect (0.66 vs.\ 0.77 with and without XPS);
\item the envelope-area effect (0.79 vs.\ 0.65 with and without it);
\item v2's polyiso assumption (89\%) and its renovation drop (0.79 $\rightarrow$ 0.40).
\end{itemize}
The same prompt revisions would have followed from the dev set alone.

\textbf{Test results.}
\begin{itemize}
\item On the 88 test buildings, v3 improves on v1 in coverage (26.1\% vs.\ 20.5\%), RMSE (69.2 vs.\ 74.1), bias (0.86 vs.\ 0.75) and CV (11.2\% vs.\ 19.2\%), but its median log error is unchanged (0.400 vs.\ 0.393).
\item None of the v1--v3 differences is significant at this size (McNemar $p=0.41$; Wilcoxon $p=0.38$).
\item v3 is significantly better than v2 (Wilcoxon $p=0.002$), and v2's regression against v1 replicates ($p=0.0002$).
\end{itemize}

\section{Feature Study}
\label{sec:features}
With v3 fixed, we varied the input features and ran each condition once (run~0), pairing it with run~0 of v3 on all 48 features. The conditions were:
\begin{itemize}
\item \emph{early}: 39 features; no LCA-derived fields and no \texttt{mass\_total}, which is the reference LCA's own inventory mass and therefore unknown at design time;
\item \emph{env}: 15 envelope-relevant features;
\item \emph{min}: 10 features (climate zone, use, project type, floor area, storeys, WWR, R-values);
\item \emph{no-lca}: all features except the 8 LCA-derived fields;
\item \emph{no-mass}: all features except \texttt{mass\_total}.
\end{itemize}

\begin{table}[t]
\centering\small
\caption{Feature study (prompt v3, run 0, 292 buildings, paired with all 48 features). Run-to-run spread of v3 itself: 2.0 points of coverage and 0.025 in median $|\log|$.}
\label{tab:features}
\begin{tabular}{lcccccc}
\toprule
Condition & Features & Coverage & Med.\ $|\log|$ & $\Delta$ Med.\ $|\log|$ & Wilcoxon $p$ & Bias \\
\midrule
all & 48 & 24.3\% & 0.461 & -- & -- & 0.86 \\
early & 39 & 21.6\% & 0.462 & $+0.001$ & 0.13 & 0.88 \\
env & 15 & 19.9\% & 0.463 & $+0.001$ & 0.33 & 0.93 \\
min & 10 & 25.3\% & 0.476 & $+0.015$ & 0.81 & 0.81 \\
no-lca & 40 & 25.3\% & 0.441 & $-0.020$ & 0.57 & 0.90 \\
no-mass & 47 & 22.6\% & 0.452 & $-0.009$ & 0.92 & 0.85 \\
\bottomrule
\end{tabular}
\end{table}

\textbf{Results} (Table~\ref{tab:features}).
\begin{itemize}
\item No feature set differs significantly from the full set, and every difference lies within the run-to-run spread of v3 itself.
\item Ten early-design inputs are as good as 48.
\item The model does not depend on \texttt{mass\_total} or the other LCA-derived fields.
\end{itemize}
Prompt content moved the median $|\log|$ error by 0.137 from v2 to v3; no feature set moved it by more than 0.020. The missing information is the assembly type, which none of the features provides.

\section{Discussion}
\textbf{Prompt design.} Three lessons transfer beyond this dataset.
\begin{enumerate}
\item A structured, step-by-step method improves consistency, but it also passes the model's component-level priors straight into the total.
\item Qualitative knowledge in a prompt is repeated in the reasoning but does not move the numbers. Quantitative, sourced anchors placed at the step where a number is chosen, with the condition under which they apply, do move them.
\item Rules that subtract (exclusions, ``only new materials'') amplify an existing low bias unless they are paired with what to include.
\end{enumerate}

\textbf{Early design and LEED v5.} An enclosure estimate within $\pm$20\% of a full LCA is reached for about a quarter of buildings, with a bias of $-14\%$. That is useful for comparing assemblies and spotting hotspots (curtain wall, HFC XPS, insulated metal panels) at schematic design. It is not reliable enough to predict a credit outcome for an individual building. The largest gains in accuracy would come from adding the assembly type to the inputs, which the proposal anticipated, rather than from further prompt engineering.

\section{Limitations}
\begin{itemize}
\item Answers within a task file share a context, so they are less independent than separate API calls.
\item The held-out split was assigned after the runs; the dev-only replication mitigates but does not replace an external test.
\item With 3 runs, the v3--v1 gain is borderline significant on all 292 buildings and not significant on the test set.
\item The feature study used one run per condition, and 7 of 9 planned leave-one-group-out ablations were stopped to limit cost, so it is a null result, not proof of irrelevance.
\item The unit is per m$^2$ of floor area, and the enclosure includes windows, doors and roofs.
\item Reference quality is uneven: there are incomplete inventories, and the tools' data vintages differ (HFC XPS).
\item Results are specific to Claude Opus~5.5 at medium effort.
\end{itemize}

\section{Conclusion}
A frontier LLM can produce enclosure A1--A3 estimates from early-stage features. Its main error is compression toward a typical value, together with blind spots for specific high-carbon materials. Error analysis on the model's own reasoning texts located those blind spots. Placing sourced quantitative anchors at the decision step reduced the bias from 0.75 to 0.86 and the run-to-run variation by a third, and raised coverage from 19.5\% to 25.0\%. Further progress depends on giving the model the assembly information that early-stage feature sets currently lack.

\begin{thebibliography}{99}\small\sloppy
\bibitem{benke2025} B.~Benke, M.~Chafart, Y.~Shen, M.~Ashtiani, S.~Carlisle, K.~Simonen. A harmonized dataset of high-resolution embodied life cycle assessment results for buildings in North America. \emph{Scientific Data} 12 (2025). \url{https://doi.org/10.1038/s41597-025-05216-0}
\bibitem{jung2026} E.~H.~Jung et al. Large language models based in-context learning for early stage building life cycle assessment. \emph{Proc.\ 59th Hawaii Int.\ Conf.\ on System Sciences (HICSS-59)}, 2026.
\bibitem{clf2025} Carbon Leadership Forum. \emph{2025 CLF North American Material Baselines} (and 2023 edition). \url{https://carbonleadershipforum.org/clf-north-american-material-baselines/}
\bibitem{alberta} Government of Alberta. \emph{PERSIST: Insulation} technical series, 2016. \url{https://www.alberta.ca/system/files/custom_downloaded_images/tr-dtseries06insul2016.pdf}
\bibitem{vermont} Vermont Public Service Department. Section R408 Insulation Embodied Carbon Calculation (draft). \url{https://publicservice.vermont.gov/}
\bibitem{oneclick} One Click LCA EPD library: aluminium-framed curtain wall systems. \url{https://materials.oneclicklca.com/}
\bibitem{arup} Arup and Saint-Gobain. \emph{The carbon footprint of facades: significance of glass}. \url{https://www.arup.com/}
\bibitem{kawneer} Kawneer (Arconic). Environmental Product Declarations for curtain wall, storefront and window systems, UL Environment, 2021.
\bibitem{pci} Precast/Prestressed Concrete Institute. Industry-average EPDs for architectural and insulated precast, 2025.
\bibitem{rdh} RDH Building Science. \emph{Embodied Carbon Analysis: Spray Foam Assemblies}, 2025; RDH \& Toronto Metropolitan University, \emph{Embodied Carbon Resource Guide}.
\bibitem{carbenmats} M.~R\"ock et al. CarbEnMats-Buildings: a global database on whole life carbon, energy and material intensity of buildings. Zenodo, 2023. \url{https://zenodo.org/records/8363895}
\end{thebibliography}

\appendix
\section{Prompts}
\label{app:prompts}
The prompts below are exported verbatim from \texttt{experiment\_config.py} by \path{report/latex/make_prompts.py}. Six non-ASCII characters in the glossary are shown as ASCII (e.g.\ CO2 for CO$_2$). Each request consists of the system prompt, the shared context, and a building block (the building's features followed by the task text).

\subsection{System prompt}
\lstinputlisting{prompts/system.txt}

\subsection{Shared context (identical in v1, v2 and v3)}
\lstinputlisting{prompts/shared_context.txt}

\subsection{Example building features (project 225, all 48 features)}
\lstinputlisting{prompts/building_225.txt}

\subsection{Task text v1 (baseline)}
\lstinputlisting{prompts/task_v1.txt}

\subsection{Task text v2}
\lstinputlisting{prompts/task_v2.txt}

\subsection{Task text v3 (final prompt)}
\label{app:v3}
The reference values are documented with sources, system boundaries and unit conversions in \path{references/envelope_intensities_sources.md}.
\lstinputlisting{prompts/task_v3.txt}

\subsection{Subagent instructions (template)}
Placeholders in braces are filled per task file by \texttt{subagent\_runs.py prepare}.
\lstinputlisting{prompts/subagent_instructions.txt}

\section{Reproduction}
All steps are scripted in the repository:
\begin{lstlisting}
python GWP-calculation/sum_gwp.py                  # ground truth (enclosure, A1-A3)
python GWP-calculation/enclosure_breakdown.py      # material breakdown + data-quality flags
python subagent_runs.py prepare --run-name 20261104_prompt-v3 --runs 3 --effort medium
#   one Claude Code subagent (Opus, effort medium) per results/<run>/subagent/tasks/*.instructions.txt
python subagent_runs.py resume  --run-name 20261104_prompt-v3   # only if interrupted
python subagent_runs.py collect --run-name 20261104_prompt-v3
python summarize_run.py 20261104_prompt-v3
python analysis/make_split.py
python analysis/compare_runs.py 20261104_baseline-innate 20261104_prompt-v2 20261104_prompt-v3
python analysis/error_analysis.py --split dev 20261104_baseline-innate 20261104_prompt-v2 20261104_prompt-v3
python analysis/feature_study.py
python report/make_figures.py
\end{lstlisting}

\end{document}
